\documentclass[11pt,a4paper]{article}
\newcommand{\intestazione}{Trimero ad anello --- il filo conduttore}
% =====================================================================
%  Preambolo comune ai documenti sul dimero (serie "che cosa abbiamo fatto")
%  Incluso con % =====================================================================
%  Preambolo comune ai documenti sul dimero (serie "che cosa abbiamo fatto")
%  Incluso con % =====================================================================
%  Preambolo comune ai documenti sul dimero (serie "che cosa abbiamo fatto")
%  Incluso con \input{preambolo_dimero} subito dopo \documentclass.
% =====================================================================
\usepackage[T1]{fontenc}
\usepackage[utf8]{inputenc}
\usepackage[main=italian,provide=*]{babel}
\usepackage{amsmath,amssymb,mathtools}
\usepackage{graphicx}
\usepackage{booktabs}
\usepackage{colortbl}
\usepackage{xcolor}
\usepackage{geometry}
\usepackage{placeins}
\usepackage{caption}
\usepackage{enumitem}
\usepackage{fancyhdr}
\usepackage{needspace}
\usepackage{tikz}
\usetikzlibrary{arrows.meta,positioning,shapes.geometric,calc,fit,backgrounds}
\usepackage{hyperref}
\hypersetup{colorlinks=true,linkcolor=blu,citecolor=blu,urlcolor=blu}

\geometry{margin=2.5cm}

% --- colori coerenti con quelli delle figure (palette Okabe-Ito)
\definecolor{blu}{HTML}{0072B2}
\definecolor{arancio}{HTML}{E69F00}
\definecolor{verdes}{HTML}{009E73}
\definecolor{rossos}{HTML}{D55E00}
\definecolor{violas}{HTML}{CC79A7}
\definecolor{grigetto}{HTML}{E8E8E8}

\captionsetup{font=small,labelfont=bf,width=0.92\textwidth}

\pagestyle{fancy}
\fancyhf{}
\fancyhead[L]{\small\itshape\intestazione}
\fancyhead[R]{\small\thepage}
\renewcommand{\headrulewidth}{0.3pt}

% --- notazione
\newcommand{\ket}[1]{\lvert #1\rangle}
\newcommand{\bra}[1]{\langle #1 \rvert}
\newcommand{\media}[1]{\langle #1 \rangle}
\newcommand{\vs}[1]{\mathbf{s}_{#1}}

% --- riquadro "in una riga": il messaggio da portare a casa
\newcommand{\messaggio}[1]{%
  \begin{center}
  \begin{minipage}{0.93\textwidth}
  \setlength{\fboxsep}{9pt}%
  \colorbox{grigetto}{\begin{minipage}{\dimexpr\textwidth-2\fboxsep}
  \small\textbf{In una riga.}\enspace #1
  \end{minipage}}
  \end{minipage}
  \end{center}
}

% --- riquadro "come lo sappiamo": la verifica numerica dietro un'affermazione
\newcommand{\verifica}[1]{%
  \begin{center}
  \begin{minipage}{0.93\textwidth}
  \small\textcolor{verdes}{\rule{2.2em}{1.1pt}}\enspace
  \textcolor{verdes}{\textbf{Come lo sappiamo.}}\enspace #1
  \end{minipage}
  \end{center}
  \medskip
}
 subito dopo \documentclass.
% =====================================================================
\usepackage[T1]{fontenc}
\usepackage[utf8]{inputenc}
\usepackage[main=italian,provide=*]{babel}
\usepackage{amsmath,amssymb,mathtools}
\usepackage{graphicx}
\usepackage{booktabs}
\usepackage{colortbl}
\usepackage{xcolor}
\usepackage{geometry}
\usepackage{placeins}
\usepackage{caption}
\usepackage{enumitem}
\usepackage{fancyhdr}
\usepackage{needspace}
\usepackage{tikz}
\usetikzlibrary{arrows.meta,positioning,shapes.geometric,calc,fit,backgrounds}
\usepackage{hyperref}
\hypersetup{colorlinks=true,linkcolor=blu,citecolor=blu,urlcolor=blu}

\geometry{margin=2.5cm}

% --- colori coerenti con quelli delle figure (palette Okabe-Ito)
\definecolor{blu}{HTML}{0072B2}
\definecolor{arancio}{HTML}{E69F00}
\definecolor{verdes}{HTML}{009E73}
\definecolor{rossos}{HTML}{D55E00}
\definecolor{violas}{HTML}{CC79A7}
\definecolor{grigetto}{HTML}{E8E8E8}

\captionsetup{font=small,labelfont=bf,width=0.92\textwidth}

\pagestyle{fancy}
\fancyhf{}
\fancyhead[L]{\small\itshape\intestazione}
\fancyhead[R]{\small\thepage}
\renewcommand{\headrulewidth}{0.3pt}

% --- notazione
\newcommand{\ket}[1]{\lvert #1\rangle}
\newcommand{\bra}[1]{\langle #1 \rvert}
\newcommand{\media}[1]{\langle #1 \rangle}
\newcommand{\vs}[1]{\mathbf{s}_{#1}}

% --- riquadro "in una riga": il messaggio da portare a casa
\newcommand{\messaggio}[1]{%
  \begin{center}
  \begin{minipage}{0.93\textwidth}
  \setlength{\fboxsep}{9pt}%
  \colorbox{grigetto}{\begin{minipage}{\dimexpr\textwidth-2\fboxsep}
  \small\textbf{In una riga.}\enspace #1
  \end{minipage}}
  \end{minipage}
  \end{center}
}

% --- riquadro "come lo sappiamo": la verifica numerica dietro un'affermazione
\newcommand{\verifica}[1]{%
  \begin{center}
  \begin{minipage}{0.93\textwidth}
  \small\textcolor{verdes}{\rule{2.2em}{1.1pt}}\enspace
  \textcolor{verdes}{\textbf{Come lo sappiamo.}}\enspace #1
  \end{minipage}
  \end{center}
  \medskip
}
 subito dopo \documentclass.
% =====================================================================
\usepackage[T1]{fontenc}
\usepackage[utf8]{inputenc}
\usepackage[main=italian,provide=*]{babel}
\usepackage{amsmath,amssymb,mathtools}
\usepackage{graphicx}
\usepackage{booktabs}
\usepackage{colortbl}
\usepackage{xcolor}
\usepackage{geometry}
\usepackage{placeins}
\usepackage{caption}
\usepackage{enumitem}
\usepackage{fancyhdr}
\usepackage{needspace}
\usepackage{tikz}
\usetikzlibrary{arrows.meta,positioning,shapes.geometric,calc,fit,backgrounds}
\usepackage{hyperref}
\hypersetup{colorlinks=true,linkcolor=blu,citecolor=blu,urlcolor=blu}

\geometry{margin=2.5cm}

% --- colori coerenti con quelli delle figure (palette Okabe-Ito)
\definecolor{blu}{HTML}{0072B2}
\definecolor{arancio}{HTML}{E69F00}
\definecolor{verdes}{HTML}{009E73}
\definecolor{rossos}{HTML}{D55E00}
\definecolor{violas}{HTML}{CC79A7}
\definecolor{grigetto}{HTML}{E8E8E8}

\captionsetup{font=small,labelfont=bf,width=0.92\textwidth}

\pagestyle{fancy}
\fancyhf{}
\fancyhead[L]{\small\itshape\intestazione}
\fancyhead[R]{\small\thepage}
\renewcommand{\headrulewidth}{0.3pt}

% --- notazione
\newcommand{\ket}[1]{\lvert #1\rangle}
\newcommand{\bra}[1]{\langle #1 \rvert}
\newcommand{\media}[1]{\langle #1 \rangle}
\newcommand{\vs}[1]{\mathbf{s}_{#1}}

% --- riquadro "in una riga": il messaggio da portare a casa
\newcommand{\messaggio}[1]{%
  \begin{center}
  \begin{minipage}{0.93\textwidth}
  \setlength{\fboxsep}{9pt}%
  \colorbox{grigetto}{\begin{minipage}{\dimexpr\textwidth-2\fboxsep}
  \small\textbf{In una riga.}\enspace #1
  \end{minipage}}
  \end{minipage}
  \end{center}
}

% --- riquadro "come lo sappiamo": la verifica numerica dietro un'affermazione
\newcommand{\verifica}[1]{%
  \begin{center}
  \begin{minipage}{0.93\textwidth}
  \small\textcolor{verdes}{\rule{2.2em}{1.1pt}}\enspace
  \textcolor{verdes}{\textbf{Come lo sappiamo.}}\enspace #1
  \end{minipage}
  \end{center}
  \medskip
}


\title{\bfseries Il trimero ad anello: che cosa abbiamo fatto, e perché\\[2pt]
\large Documento 0 --- il filo conduttore}
\author{Samuele Schiavi \\ \small Tesi triennale in Fisica, Università di Parma
--- Relatore: Prof.\ Alessandro Chiesa}
\date{}

\begin{document}
\maketitle
\thispagestyle{fancy}

\begin{abstract}
\noindent
Questa serie racconta l'estensione a tre spin, in topologia ad anello, degli
strumenti già validati sul dimero (serie \emph{Dimero di spin}, Documenti 0--4):
stesso ordine di lavoro, stessa disciplina di verifica, stessi quattro blocchi.
Non è una ripetizione: tre spin accoppiati ciclicamente introducono
\textbf{frustrazione geometrica} (assente nel dimero) e complicazioni tecniche
concrete --- un ansatz a sei parametri invece di tre, un secondo livello di
approssimazione di Trotter, un circuito dei correlatori a quattro qubit --- che
mettono alla prova se il criterio di progetto del dimero regge quando il
sistema cresce. Ogni numero riportato è stato ricalcolato eseguendo il codice
di questo pacchetto, non ripreso da sessioni precedenti.
\end{abstract}

\section{Perché tre spin, e perché ad anello}

Tre spin $1/2$ danno uno spazio di Hilbert di dimensione $8$ --- ancora
piccolo abbastanza da diagonalizzare esattamente, ma abbastanza grande da
introdurre una fisica che il dimero non può mostrare per costruzione: un
numero \emph{dispari} di spin $1/2$ accoppiati antiferromagneticamente in un
ciclo chiuso non può soddisfare simultaneamente tutti i legami --- è
\textbf{frustrazione geometrica}, la stessa fisica che rende interessanti i
triangoli di spin nei materiali molecolari reali.

Fra le due topologie possibili a $N=3$ (anello chiuso, o catena aperta), il
lavoro qui documentato usa l'\textbf{anello}: genuinamente frustrato, ansatz
più semplice (un solo giro di blocchi basta, si vedrà nel Documento 2), meno
correlatori indipendenti da verificare per via degli zeri di simmetria. La
catena aperta resta un filone possibile, non affrontato in questa serie.

\messaggio{Il trimero ad anello non ripete il dimero su un sistema più
grande: aggiunge un ingrediente fisico --- la frustrazione --- che il dimero,
per costruzione (due soli spin, un solo legame), non può contenere.}

\section{I quattro passi, e come si incastrano}

Stessa struttura del dimero, stesso ordine di dipendenze.

\begin{figure}[htbp]
\centering
\resizebox{\textwidth}{!}{%
\begin{tikzpicture}[
  scatola/.style={draw,rounded corners=3pt,minimum height=1.45cm,
                  minimum width=3.0cm,align=center,font=\small,thick},
  freccia/.style={-{Stealth[length=3mm]},thick,gray}]
\node[scatola,draw=blu,fill=blu!8]      (A) at (0,0)    {\textbf{1. Sistema}\\[1pt]
   spettro esatto\\ frustrazione, DM};
\node[scatola,draw=verdes,fill=verdes!8](B) at (5.0,0)  {\textbf{2. VQE}\\[1pt]
   stato fondamentale\\ RBS contro $W$};
\node[scatola,draw=arancio,fill=arancio!10](C) at (10.0,0){\textbf{3. Dinamica}\\[1pt]
   $e^{-iHt}$, due livelli\\ di Trotter annidati};
\node[scatola,draw=rossos,fill=rossos!8](D) at (15.0,0) {\textbf{4. Correlatori}\\[1pt]
   ancilla, 81 combinazioni\\ 4 qubit totali};
\draw[freccia] (A) -- node[above,font=\scriptsize,text=black]{riferimento}
                        node[below,font=\scriptsize,text=black]{esatto} (B);
\draw[freccia] (B) -- node[above,font=\scriptsize,text=black]{ansatz}
                        node[below,font=\scriptsize,text=black]{canonico} (C);
\draw[freccia] (C) -- node[above,font=\scriptsize,text=black]{$U(t)$}
                        node[below,font=\scriptsize,text=black]{nel circuito} (D);
\end{tikzpicture}}
\caption{La catena di dipendenze, identica nella forma a quella del dimero. Lo
spettro esatto del passo 1 è il metro di giudizio per il VQE; l'ansatz
canonico scelto nel passo 2 (non uno stato) è riusato, ri-ottimizzato punto
per punto, nei passi 3 e 4.}
\end{figure}
\FloatBarrier

\section{L'Hamiltoniana}

Trimero isoscele: legame $J$ fra i siti $(1,2)$, legami $J'$ fra $(2,3)$ e
$(3,1)$ --- l'equilatero frustrato è il caso particolare $J'=J$. Con
$J'\ne0$ il ciclo si chiude: è questo legame in più, assente nella catena
aperta, a rendere il sistema genuinamente frustrato.
\begin{equation}
H = \underbrace{J\,\vec\sigma_1\!\cdot\!\vec\sigma_2 +
J'\,(\vec\sigma_2\!\cdot\!\vec\sigma_3+\vec\sigma_3\!\cdot\!\vec\sigma_1)}_{\text{scambio isotropo, frustrato}}
+ \underbrace{b\,(Z_1{+}Z_2{+}Z_3)}_{\text{campo}}
+ \underbrace{D\sum_{\langle i,j\rangle} (X_iZ_j-Z_iX_j)}_{\text{termine DM}}.
\label{eq:H}
\end{equation}
Ogni spin $1/2$ resta un qubit, come sul dimero: nessun costo di
rappresentazione aggiuntivo nel passare da due a tre siti. Il rapporto
$J'/J=0.4$ è fissato in tutta la serie (lo stesso usato nel lavoro
preliminare sulla teoria esatta); $J=1$ fissa l'unità di energia.

\section{Il termine DM: due opzioni, una scelta}

A differenza del dimero (dove il DM ha una forma unica), con tre legami ci
sono più modi di accendere un termine DM, e non tutti aprono l'incrocio
allo stesso modo:
\begin{itemize}[leftmargin=1.5em,itemsep=2pt]
\item \textbf{Opzione A}: DM nullo sul legame base, segno opposto sui due
legami laterali. Rispetta una simmetria residua dell'Hamiltoniana
($S_{12}^2$, il quadrato dello spin della coppia base) --- e per questo
\textbf{non apre mai l'incrocio fondamentale}, qualunque sia la sua
intensità.
\item \textbf{Opzione B}: DM su tutti e tre i legami, stesso segno,
proporzionale a $J$ o $J'$. Rompe quella simmetria residua e \textbf{apre
l'incrocio davvero}.
\end{itemize}

\paragraph{Verifica sistematica dei segni.} Non una scelta fra due
candidati equivalenti: costruita esplicitamente la matrice di permutazione
$P_{12}$ (scambio dei siti $1,2$) e verificato
$\max|P_{12}H_\text{DM}P_{12}^\dagger-H_\text{DM}|$ per tutte le
combinazioni di segno di $(D_{12},D_{23},D_{31})\propto(r{\cdot}J,r{\cdot}J',r{\cdot}J')$,
a $r=0.15$ (dettaglio completo in \texttt{analisi\_dm\_trimero\_anello.tex},
riprodotto e verificato indipendentemente qui):

\begin{table}[htbp]
\centering
\small
\renewcommand{\arraystretch}{1.25}
\begin{tabular}{@{}cccc@{}}
\toprule
\rowcolor{grigetto}
$D_{12}$ & $D_{23}$ & $D_{31}$ & $\max|[P_{12},H_\text{DM}]|$ \\
\midrule
$0$ & $+$ & $+$ & $0.2400$ \\
$0$ & $+$ & $-$ & $\mathbf{0.0000}$ \quad (simmetrica --- Opzione A) \\
$0$ & $-$ & $+$ & $\mathbf{0.0000}$ \quad (simmetrica) \\
$0$ & $-$ & $-$ & $0.2400$ \\
$+$ & $+$ & $+$ & $0.4200$ \quad (Opzione B) \\
$+$ & $+$ & $-$ & $0.3000$ \\
$+$ & $-$ & $+$ & $0.3000$ \\
$+$ & $-$ & $-$ & $0.4200$ \\
$-$ & $+$ & $+$ & $0.4200$ \\
$-$ & $+$ & $-$ & $0.3000$ \\
$-$ & $-$ & $+$ & $0.3000$ \\
$-$ & $-$ & $-$ & $0.4200$ \\
\bottomrule
\end{tabular}
\caption{Tutte e 12 le combinazioni di segno testate. Solo
$(0,+,-)$/$(0,-,+)$ rispettano la simmetria; l'Opzione B ($0.42$) è, fra
tutte, quella che la rompe di più --- non una rottura qualsiasi, la
massima fra i candidati naturali (proporzionali a $J,J'$).}
\end{table}
\FloatBarrier

\begin{table}[htbp]
\centering
\renewcommand{\arraystretch}{1.4}
\begin{tabular}{@{}p{2.3cm}p{6.3cm}p{6.3cm}@{}}
\toprule
\rowcolor{grigetto}
& \textbf{Opzione A} (simmetrica) & \textbf{Opzione B} (proposta del relatore) \\
\midrule
forma & $D'\big[(X_2Z_3{-}Z_2X_3)-(X_3Z_1{-}Z_3X_1)\big]$ &
$D(X_1Z_2{-}Z_1X_2)+0.4D(X_2Z_3{-}Z_2X_3)+0.4D(X_3Z_1{-}Z_3X_1)$ \\
simmetria $P_{12}$ & rispettata esattamente & rotta \\
$[H,S_{12}^2]$ & $=0$ esatto ($\sim10^{-16}$) & $\neq0$ ($\sim1.4$ a
$D{=}0.15$) \\
struttura residua & blocco $2{\times}2$ (C) $+$ blocco $6{\times}6$
(A+B misti) & nessuna, $8\times8$ pieno \\
\bottomrule
\end{tabular}
\caption{Le due opzioni rimaste dopo la verifica sistematica. Qui $D$ è il
parametro di intensità complessivo (chiamato $r$ nel Documento 1); i
coefficienti $0.4$ sull'Opzione B vengono da $J'/J=0.4$, coincidono
numericamente con $J'$ solo perché $J=1$.}
\end{table}
\FloatBarrier

Il Documento 1 verifica entrambe le affermazioni direttamente. La scelta
operativa per il resto della serie (VQE, dinamica, correlatori) è
l'\textbf{Opzione B} --- l'unica che riproduce, per il trimero, lo stesso
ruolo che il DM ha nel dimero.

\messaggio{Con tre spin non basta dire ``il DM rompe la simmetria'': bisogna
specificare quale forma di DM, perché non tutte le forme possibili lo fanno.
È una domanda che il dimero, con un solo legame, non poteva nemmeno porre.}

\section{I due punti di lavoro usati}

\begin{table}[htbp]
\centering
\footnotesize
\renewcommand{\arraystretch}{1.35}
\begin{tabular}{@{}p{2.5cm}p{3.6cm}p{2.3cm}p{5.8cm}@{}}
\toprule
\rowcolor{grigetto}
\textbf{Nome} & \textbf{Parametri} & \textbf{Dove è usato} & \textbf{Perché quello} \\
\midrule
Scansione in campo & $D=0$, Opzioni A/B & Documento 1 & serve vedere l'incrocio e il ruolo del DM \\
Scenario A & $b/J=b_c=2.4$, $D/J=0.15$ (Opzione B) & Documenti 2, 4 & campo critico esatto, DM ben acceso: la condizione giusta per un confronto fra ansatz interessante \\
Punto $R_0$ & $b/J=0.05$, $D/J=1.93$ & Documento 3 & dinamica non monocromatica; \emph{dichiarato provvisorio}, non ancora derivato in modo principiato \\
\bottomrule
\end{tabular}
\caption{I punti di lavoro. Come sul dimero, il punto della dinamica e quello
del VQE/correlatori sono diversi apposta: la condizione per una dinamica
ricca (più frequenze visibili) non coincide con quella per un fondamentale
non banale. A differenza del dimero, qui l'ansatz canonico (Documento 2) è
validato in un punto singolo, non su uno sweep in campo --- si vedrà nel
Documento 2 perché uno sweep con questo ansatz non è, allo stato attuale,
attendibile su tutto il range.}
\end{table}
\FloatBarrier

\section{Come sono state ottenute le cifre di questi documenti}

Stessi tre criteri del dimero, applicati identici.

\begin{enumerate}[leftmargin=1.5em,itemsep=3pt]
\item \textbf{Nessun numero preso per buono dalla teoria.} Ogni valore
proviene dall'esecuzione del codice di questo pacchetto.
\item \textbf{Due metodi indipendenti dove è possibile --- e dichiarato
quando non lo è.} Il circuito dei correlatori è verificato contro
l'esponenziale di matrice diretto della piena Hamiltoniana: derivazione
indipendente, nessun codice condiviso. La ``matrice compatta'' di Trotter,
invece, è costruita estraendo l'operatore Qiskit dello stesso circuito a
un passo (\texttt{Operator(qc)}) ed elevandolo a potenza --- non una
derivazione indipendente, ma un controllo di coerenza interna (conferma
che l'elevamento a potenza riproduce correttamente $N$ ripetizioni del
circuito), utile ma di natura diversa: può scoprire un bug di
implementazione, non un errore nel modello fisico.
\item \textbf{Le discrepanze si riportano, non si smussano.} Dove emerge un
comportamento diverso dal dimero --- e succede più di una volta in questa
serie --- è dichiarato esplicitamente, non nascosto nella scala di un
grafico.
\end{enumerate}

\verifica{Come prima applicazione del criterio: lo spettro numerico ottenuto
diagonalizzando la matrice $8\times8$ coincide con le energie analitiche dei
tre multipletti di Kambe (blocchi A, B, C) a precisione macchina su tutta la
griglia in $b$ usata nel Documento 1.}

\section{Che cosa non c'è in questa serie}

\begin{itemize}[leftmargin=1.5em,itemsep=2pt]
\item \textbf{Il rumore.} Come sul dimero, tutto qui è a simulazione ideale.
Il sistema come sistema aperto (rumore di gate e di lettura) è trattato in un
pacchetto separato, che riusa esattamente l'ansatz e i punti di lavoro
stabiliti qui.
\item \textbf{La catena aperta.} Topologia alternativa a $N=3$, con teoria
esatta già disponibile ma non ripresa in questa serie.
\end{itemize}

\end{document}
